\documentclass[letterpaper,11pt]{article}

\usepackage[T1]{fontenc}
\usepackage[utf8]{inputenc}
\usepackage{mathptmx}            % Times
\usepackage[left=0.75in,right=0.75in,top=0.5in,bottom=0.5in]{geometry}
\usepackage{titlesec}
\usepackage{enumitem}
\usepackage[hidelinks]{hyperref}

\pagestyle{empty}
\setlength{\parindent}{0pt}
\setlength{\parskip}{0pt}

% Section headings: normal weight, rule underneath
\titleformat{\section}{\large}{}{0pt}{}[\vspace{-12pt}\rule{\textwidth}{0.6pt}]
\titlespacing*{\section}{0pt}{6pt}{3pt}

% Bullets
\setlist[itemize]{leftmargin=0.22in,label=\textbullet,itemsep=0pt,parsep=0pt,topsep=1pt,partopsep=0pt}

% Role / project heading: \role{Title}{Org}{Dates}{separator}
\newcommand{\role}[4]{\textbf{#1} #4 \textbf{\textit{#2}} (#3)}

\begin{document}

%-------------------- Header --------------------
\begin{center}
  {\fontsize{17}{20}\selectfont\textbf{Eric Goossen}}\\[2pt]
  \small
  \href{mailto:ericjgoossen@gmail.com}{ericjgoossen@gmail.com} $|$
  \href{https://github.com/EricJGoossen}{github.com/EricJGoossen} $|$
  \href{https://linkedin.com/in/ericjgoossen}{linkedin.com/in/ericjgoossen}
\end{center}
\vspace{-8pt}

\small

%-------------------- Education --------------------
\section{Education}
\textbf{University of Washington, Seattle,} Computer Engineering BS (Jun 2028 expected), GPA: 3.88

%-------------------- Skills --------------------
\section{Technical Skills}
\textbf{Languages:} C++, Python, Rust \;$|$\; \textbf{Tools:} ROS2, Isaac Sim, MuJoCo, Pinocchio, Docker, Git, Linux\\
\textbf{Robotics:} Kinematics \& Dynamics, Motion Planning, Impedance Control, LQR, Kalman Filtering, SLAM\\
\textbf{Hardware:} Jetson Orin/Thor, Franka, OpenArm, CAN-FD, LiDAR, ESP32, oscilloscope

%-------------------- Experience --------------------
\section{Experience}
\role{Undergraduate Research Intern}{UW Human Centered Robotics Lab}{Apr 2026 -- Present}{--}
\begin{itemize}
  \item Co-authored research testing whether human feedback helps LLM-written robot code beat vision-language-action models
  \item Showed human feedback lifts Code-as-Policies 1.75x over base VLAs, solving tasks VLAs failed completely
  \item Sole author of the paper's open-source VLA eval framework, deploying chunked actions to a Franka via DROID
  \item Built one policy interface over WebSocket, ZMQ, and HTTP so any VLA can be evaluated with no custom code
  \item Benchmarked $\pi$0, $\pi$0.5, GR00T N1.7, and MolmoAct2 on 12 real-world manipulation tasks with automated logging
  \item Built Oculus teleop demo collection for each task, then fine-tuned $\pi$0, $\pi$0.5, and MolmoAct2 on the demos
\end{itemize}

\role{Undergraduate Research Intern}{UW Personal Robotics Lab}{Oct 2025 -- Present}{--}
\begin{itemize}
  \item Lead 4-person undergrad team building the software stack for an NVIDIA-granted bimanual holonomic robot
  \item Own base navigation, localization, and arm control, built in MuJoCo and Gazebo, then deployed to hardware
  \item Robot navigates point to point autonomously; both arms move in 3D under autonomous and teleop control
  \item Extended MuJoCo planner to dual-arm with cross-arm collision checking, 125Hz retiming, and parallel execution
  \item Fixed and upstreamed open-source inverse kinematics solver bug, from failing to 50/50 on library test cases
  \item Prototyped three impedance controllers; shipped 500Hz gravity comp (Pinocchio RNEA) feeding motor controllers
  \item Built sim-to-real tuning pipeline ($\sim$5 mm accuracy) and 8-stage operator-gated hardware validation suite
  \item Integrated ROS2 Nav2 with LiDAR SLAM for autonomous navigation in unknown, dynamic environments
  \item Built WebXR/Meta Quest teleop from scratch for sim and real hardware with Cartesian speed limits
\end{itemize}

\role{Control Systems Programmer}{Advanced Robotics at UW}{Oct 2025 -- Present}{-}
\begin{itemize}
  \item Implemented Kalman-filter odometry for sensor fusion, with data-driven tuning on physical competition robots
  \item Led path-tracking rewrite on embedded RTOS, replacing pure pursuit with LQR that models robot dynamics
  \item Built sensor-agnostic mocap pipeline to characterize lidar-inertial localization, from guided data collection to fitting
  \item Fit Gauss-Markov error model validated by Monte Carlo on held-out data: 5.1 cm horizontal, 0.6 cm vertical, $<$0.5$^\circ$
\end{itemize}

\role{Student Advocate}{Swerve Robotics}{Jan 2024 -- Aug 2025}{-}
\begin{itemize}
  \item Led campaign to raise \$50,000+ for a Botswana robotics advocacy trip; built curriculum and taught 50+ students
  \item Presented on Botswana national TV; worked with U.S. Embassy to lobby ministers for reduced STEM kit tariffs
\end{itemize}

%-------------------- Projects --------------------
\section{Projects}
\role{Achilles}{Rigid-Body Simulator}{2026 -- Present}{--}
\begin{itemize}
  \item Developing C++ rigid-body simulator that batches robot instances across SIMD lanes to speed up RL training
  \item Cut dynamics from O(n\textsuperscript{3}) to O(n) with an Articulated Body Algorithm engine, validated to 1e-4 rad
  \item Replaced Eigen with custom xsimd spatial algebra and a depth-sorted layout that keeps tree dynamics SIMD-safe
  \item Cut max energy error by over half using substepped Euler, Verlet, and implicit midpoint integrators
  \item Architected compile-time, data-oriented engine so new algorithms need no engine changes; CI enforces death tests
  \item Building branchless SIMD broad-phase collision that tests each body against a full lane group, packing hits into bitmasks
  \item Unified spheres, capsules, and boxes as one sphere-swept-box shape, with MuJoCo-style margins and parent/child filtering
  \item 1.65M robot-steps/s single-threaded (4$\times$ via batching), 2$\times$ MuJoCo's per-robot rate, matching its dynamics to $\sim$1e-8 error
\end{itemize}

\role{Live View}{Computer Vision Pre-Processing}{2024 -- 2025}{--}
\begin{itemize}
  \item Converted camera frames into braille-pixel representations engineered to fit FTC bitmap rules
  \item Applied image dithering and color encoding to compress frames while staying within FTC bandwidth restrictions
\end{itemize}

\role{SKIDADDLE}{Trajectory Planning Library}{2023 -- 2025}{--}
\begin{itemize}
  \item Published real-time B\'ezier path traversal library for FTC, with a derived optimal-path and convergence model
\end{itemize}

%-------------------- Publications --------------------
\section{Publications}
\textbf{``A Few Words Go a Long Way,''} D. Chen, A. R. Jain, \textbf{E. Goossen}, E. Romig, M. Murray, N. Walker, M. Cakmak\\
(\textit{in preparation for IEEE RA-L}), 2026.

\end{document}
